\begin{aiblock}
\ai{Ported from a standalone draft (\texttt{temp/simplified.tex}, since removed) and translated into this project's conventions: body indices $1$ (Frame+Rider), $2a$/$2b$ (Rear/Front wheel), $3b$ (Crankset, cassette folded into $2a$ per \texttt{wheel.tex}); frame $i$ (fixed) / $0$ (ground-local); $\theta$ for ground/frame tilt, $\phi$ for body spin (was $\theta_1$ and $\theta_{2,3,4}$ respectively -- the original's road-gradient $\theta_1$ collided with this project's spin-angle symbol once translated, hence the split). Point labels ($A$--$K$, $a$) are local to this derivation and are \emph{not} the same points as in 2D -- Systems (e.g. $C$ here is the rear wheel's shifted ground contact, not the drivetrain bearing point in \texttt{frame.tex}). Air density $\rho$ in \S Analysis is unrelated to the road curvature radius $\rho(x)$ in 2D -- Assembly; both are standard symbols in their own contexts. $M$ is the total lumped system mass (frame+rider+wheels+drivetrain combined, per the point-mass assumption below) -- distinct from a per-body $m_k$. Equation numbers restart at each of Bike, Wheels, Drivetrain, Assembly, Analysis, per the document's \texttt{numberwithin\{equation\}\{subsection\}}.}

\ai{\textbf{Axis-direction note:} within this derivation only, $\hat z_i$ and $\hat z_0$ point \emph{downward} (along gravity), the opposite of 2D's global-assumption convention (where $\hat z_i$ points up, opposite gravity). This is inherited unchanged from the original draft and kept as-is rather than resigning the whole derivation, since every equation below is internally consistent under it -- e.g.\ $\vv W_1=+Mg\hat z_i$ and the Limits condition ${}^{z_0}R_{2a/0}\le0$ ("no wheel lift") only both hold together if $z$ points down. If this gets confusing in practice, the fix is a single global sign flip on every $z$-component here, not a symbol rename.}
\end{aiblock}

\begin{aiblock}
\subsection{Bike}
The equation of motion is written for the total bike system. Internal forces between subsystems cancel and do not appear.

\subsubsubsection{Assumptions}
\begin{itemize}
    \item Rigid frame --- no compliance, no rider motion relative to frame
    \item Point mass at fixed height $h$, total mass $M$ --- rider and frame CG collocated and invariant
    \item Motion confined to the $x$--$z$ plane --- straight-line travel, no yaw or roll
    \item Constant road gradient $\theta$ --- no curvature or undulation
    \item Two contact points only --- front and rear tyre, no additional supports
    \item Aerodynamic lift neglected --- drag and pitching moment only
    \item Wind in the $x$--$z$ plane --- no crosswind
    \item Rolling resistance as contact resultant shift (Option A) --- no separate resistance moment
    \item Both tyres share the same $c_{2a,rr}=c_{2b,rr}$ and $r_{2a}=r_{2b}$
\end{itemize}

\subsubsubsection{Forces}
Weight acts at the system centre of gravity $a$ in the fixed frame:
\begin{equation}
    \vv{W}_1 = Mg\,\hat{z}_i
\end{equation}

Under load, tire deformation shifts the contact resultant forward of the geometric contact by $l_{2a,rr} = c_{2a,rr}r_{2a}$ (rear), $l_{2b,rr} = c_{2b,rr}r_{2b}$ (front). Points $C$ and $D$ denote these shifted contact locations. The ground reactions at $C$ and $D$ are decomposed in the ground-local frame $(\hat{x}_0, \hat{z}_0)$:
\begin{gather}
    \vv{R}_{2a/0} = {}^{x_0}R_{2a/0}\,\hat{x}_0 + {}^{z_0}R_{2a/0}\,\hat{z}_0 \\
    \vv{R}_{2b/0} = {}^{x_0}R_{2b/0}\,\hat{x}_0 + {}^{z_0}R_{2b/0}\,\hat{z}_0
\end{gather}

Aerodynamic drag and pitching moment act on the frame, neglecting lift:
\begin{gather}
    \vv{A}_{1,f} = -A_{1,x}\,\hat{x}_0 \\
    \vv{A}_{1,m} = -A_{1,m}\,\hat{y}
\end{gather}

\subsubsubsection{Geometry}
Horizontal distances $l_C$, $l_D$ are measured from the system CG $a$ to the geometric wheel contact points. The rolling resistance offsets $l_{2a,rr}$, $l_{2b,rr}$ shift the contact resultant forward; $\sigma$ ensures this shift is always in the direction of motion:
\begin{gather}
    L = l_C + l_D \qquad \sigma = \sgn(\dot{x}^{(1)}) \\
    l_C' = l_C - l_{2a,rr}\sigma \qquad l_D' = l_D + l_{2b,rr}\sigma
\end{gather}
The position vectors from $a$ to each force application point, expressed in the ground-local frame:
\begin{gather}
    \vv{aC} = (-l_C + l_{2a,rr}\sigma)\,\hat{x}_0 + h\,\hat{z}_0 \\
    \vv{aD} = (+l_D + l_{2b,rr}\sigma)\,\hat{x}_0 + h\,\hat{z}_0 \\
    \vv{aA} = -l_C\,\hat{x}_0 + (h - r_{2a})\,\hat{z}_0
\end{gather}

\subsubsubsection{Equations of motion}
The moment equation uses $\dot{\vv{H}}_a$, the rate of change of angular momentum about $a$, rather than a single body inertia, because the spinning wheels and crankset each contribute angular momentum independently:
\begin{align}
    M\,\ddot{\vv{x}}^{(1)} &= \sum \vv{F}_1 = \vv{W}_1 + \vv{R}_{2a/0} + \vv{R}_{2b/0} + \vv{A}_{1,f} \\[6pt]
    \dot{\vv{H}}_a &= \sum \vv{M}_1(a) = \vv{aC}\times\vv{R}_{2a/0} + \vv{aD}\times\vv{R}_{2b/0} + \vv{aA}\times\vv{A}_{1,f} + \vv{A}_{1,m}
\end{align}

\subsubsubsection{Translation}
Weight is projected onto the ground-local frame using the gradient angle $\theta$. The vertical equation enforces the constraint that the bike remains on the road:
\begin{align}
    M\,\ddot{x}^{(1)} &= -Mg\sin\theta + {}^{x_0}R_{2a/0} + {}^{x_0}R_{2b/0} - A_{1,x} \\
    0                 &= Mg\cos\theta + {}^{z_0}R_{2a/0} + {}^{z_0}R_{2b/0}
\end{align}

\subsubsubsection{Rotation}
Since all subsystems translate together, the transport terms in $\dot{H}_a$ vanish and only the spin inertias remain. The frame's own rotation is zero ($\ddot\theta = 0$, since $\theta$ is held constant here), so:
\begin{align}
    \dot{H}_y &= I_{2a}\ddot{\phi}_{2a} + I_{2b}\ddot{\phi}_{2b} + I_{3b}\ddot{\phi}_{3b} \\
    \dot{H}_y &= h\left({}^{x_0}R_{2a/0} + {}^{x_0}R_{2b/0}\right) + l_C'\,{}^{z_0}R_{2a/0} - l_D'\,{}^{z_0}R_{2b/0}
               - \left(h - r_{2a}\right)A_{1,x} - A_{1,m}
\end{align}

\subsection{Wheels}
Each wheel is an independent rigid body. Moments are taken about the axle to eliminate the axle reaction, which passes through it and produces no moment.

\subsubsubsection{Assumptions}
\begin{itemize}
    \item Each wheel is a rigid body with uniform mass distribution --- CG at the axle
    \item Axle fixed relative to the frame --- no suspension travel
    \item Wheel spin axis is $\hat{y}$ --- no camber or toe
    \item Contact patch is a point --- no contact patch width or length effects
    \item No lateral forces --- pure longitudinal contact resultant in the $x$--$z$ plane
    \item Brake force acts at a single effective radius ($r_{2a,bk}$ rear, $r_{2b,bk}$ front) --- caliper geometry collapsed
    \item No wheel flex or tyre sidewall compliance
\end{itemize}

\subsubsubsection{Geometry}
Both wheels share the same radius ($r_{2a}=r_{2b}$) and rolling resistance coefficient ($c_{2a,rr}=c_{2b,rr}$), so the axle-to-contact vector has the same form for both. The chain and brake moment arms are specific to each wheel:
\begin{equation}
    |\vv{AF}| = r_{2a,cg} \qquad |\vv{AJ}| = r_{2a,bk} \qquad |\vv{BK}| = r_{2b,bk}
\end{equation}
The axle-to-contact resultant vector, expressed in the ground-local frame:
\begin{gather}
    \vv{AC} = l_{2a,rr}\sigma\,\hat{x}_0 + r_{2a}\,\hat{z}_0 \\
    \vv{BD} = l_{2b,rr}\sigma\,\hat{x}_0 + r_{2b}\,\hat{z}_0
\end{gather}

\subsubsubsection{Forces}
The chain tension acts only on the rear wheel at the cog contact point $F$. Brake forces act tangentially at the effective pad radius on whichever wheel is braked:
\begin{gather}
    \vv{T}_{2a/3b} = T_{2a/3b}\,\hat{e}_{span} \\
    \vv{B}_{2a/1} = B_{2a/1}\,\hat{e}_{t,bk}, \qquad \vv{B}_{2b/1} = B_{2b/1}\,\hat{e}_{t,bk}
\end{gather}

\subsubsubsection{Equations of motion}
Expanding the cross products with the geometry above, noting that $\vv{R}_{2a/1}$ and $\vv{W}_{2a}$ both act through the axle and vanish (likewise $\vv{R}_{2b/1}$, $\vv{W}_{2b}$):
\begin{align}
    I_{2a}\ddot{\phi}_{2a} &= \sum \vv{M}_{2a}(A)\cdot\hat{y}
        = \bigl[\vv{AC}\times\vv{R}_{2a/0} + \vv{AF}\times\vv{T}_{2a/3b}
          + \vv{AJ}\times\vv{B}_{2a/1}\bigr]\cdot\hat{y} \\[6pt]
    I_{2b}\ddot{\phi}_{2b} &= \sum \vv{M}_{2b}(B)\cdot\hat{y}
        = \bigl[\vv{BD}\times\vv{R}_{2b/0}
          + \vv{BK}\times\vv{B}_{2b/1}\bigr]\cdot\hat{y}
\end{align}

\subsubsubsection{Rear wheel}
The traction force ${}^{x_0}R_{2a/0}$ drives the wheel; the normal load ${}^{z_0}R_{2a/0}$ contributes a rolling resistance moment through the offset $l_{2a,rr}\sigma$; chain tension resists spin; braking resists spin:
\begin{equation}
    I_{2a}\ddot{\phi}_{2a} = r_{2a}\,{}^{x_0}R_{2a/0} - l_{2a,rr}\sigma\,{}^{z_0}R_{2a/0}
        - r_{2a,cg}\,T_{2a/3b} + \sigma\,r_{2a,bk}\,B_{2a/1}
\end{equation}

\subsubsubsection{Front wheel}
The front wheel carries no chain force. Only the contact resultant and the brake act:
\begin{equation}
    I_{2b}\ddot{\phi}_{2b} = r_{2b}\,{}^{x_0}R_{2b/0} - l_{2b,rr}\sigma\,{}^{z_0}R_{2b/0}
        + \sigma\,r_{2b,bk}\,B_{2b/1}
\end{equation}

\subsection{Drivetrain}
The drivetrain subsystem comprises the crankset: chainring, cranks, and pedals. The cassette is treated as part of the rear wheel (matches Rear (2a) already carrying the cassette reaction, per \texttt{wheel.tex}). A fixed gear ratio $r_{2a,cg}/r_{3b,cr}$ is assumed throughout.

\subsubsubsection{Assumptions}
\begin{itemize}
    \item Fixed gear ratio --- no derailleur shifting, single chainring and cog
    \item No drivetrain efficiency loss --- chain transmits torque without friction
    \item Chain mass neglected --- chain inertia not modelled
    \item Crankset is a rigid body with CG at the BB axle $E$
    \item Pedal force acts tangentially at the crank tip --- no radial component
    \item Left and right cranks combined as a single effective torque
    \item No freewheel modelled explicitly --- $T_{2a/3b} \ge 0$ enforced as a post-check
\end{itemize}

\subsubsubsection{Geometry}
$E$ is the bottom bracket axle, the rotation centre of the crankset. The chainring engages the chain at $G$; the rider applies force at the pedal spindles $H_L$, $H_R$. The relevant moment arms are:
\begin{equation}
    |\vv{EG}| = r_{3b,cr} \qquad |\vv{EH_L}| = |\vv{EH_R}| = r_{1,pedal}
\end{equation}

\subsubsubsection{Forces}
The chain pulls the chainring at $G$ with the same tension it pulls the cog, by Newton's third law. The rider applies pedal forces tangentially at each crank tip:
\begin{gather}
    \vv{T}_{3b/2a} = -\vv{T}_{2a/3b} \\
    \vv{P}_{3b/1,L} = P_{3b/1,L}\,\hat{e}_{t,L}, \qquad \vv{P}_{3b/1,R} = P_{3b/1,R}\,\hat{e}_{t,R}
\end{gather}

\subsubsubsection{Equations of motion}
Moments about $E$ eliminate the BB axle reaction and crankset weight:
\begin{equation}
    I_{3b}\ddot{\phi}_{3b} = \sum \vv{M}_{3b}(E)\cdot\hat{y}
        = \bigl[\vv{EG}\times\vv{T}_{3b/2a}
          + \vv{EH_L}\times\vv{P}_{3b/1,L}
          + \vv{EH_R}\times\vv{P}_{3b/1,R}\bigr]\cdot\hat{y}
\end{equation}

\subsubsubsection{Crankset}
Pedal force drives the crankset; chain tension resists it:
\begin{equation}
    I_{3b}\ddot{\phi}_{3b} = r_{3b,cr}\,T_{2a/3b} - r_{1,pedal}\left(P_{3b/1,L} + P_{3b/1,R}\right)
\end{equation}

\subsection{Assembly}
The four subsystems are coupled through three kinematic constraints: no-slip at each contact and the fixed chain ratio. All angular accelerations reduce to a single unknown $\ddot{x}^{(1)}$. The solve order has no algebraic loop: the EOM gives $\ddot{x}^{(1)}$ first, then reactions are post-processed in sequence.

\subsubsubsection{Assumptions}
\begin{itemize}
    \item No-slip at both contacts --- wheels roll without sliding
    \item Kinematic constraint is holonomic --- $\dot{x}^{(1)} = -r_{2a}\dot{\phi}_{2a} = -r_{2b}\dot{\phi}_{2b}$ at all times
    \item Chain does not slip on sprockets --- gear ratio is exact
    \item All assumptions from Bike, Wheels, and Drivetrain carry forward
\end{itemize}

\subsubsubsection{Kinematics}
No-slip at each contact relates wheel angular acceleration to frame linear acceleration. The chain ratio propagates this to the crankset:
\begin{align}
    \ddot{\theta} &= 0 \\
    \ddot{\phi}_{2a} = \ddot{\phi}_{2b} &= -\frac{\ddot{x}^{(1)}}{r_{2a}} \\
    \ddot{\phi}_{3b} &= -\frac{r_{2a,cg}}{r_{3b,cr}\,r_{2a}}\,\ddot{x}^{(1)}
\end{align}

\subsubsubsection{Chain tension}
Substituting kinematics into the crankset equation and solving for $T_{2a/3b}$. The first term captures crankset inertia; the second is the pedal input referred to the chain:
\begin{equation}
    \boxed{T_{2a/3b} = -\frac{I_{3b}\,r_{2a,cg}}{r_{3b,cr}^2\,r_{2a}}\,\ddot{x}^{(1)}
           + \frac{r_{1,pedal}}{r_{3b,cr}}\left(P_{3b/1,L} + P_{3b/1,R}\right)}
\end{equation}

\subsubsubsection{Traction}
Substituting kinematics into each wheel equation and solving for the longitudinal ground reaction. These are the horizontal forces the ground must supply; they are checked against the friction limit below:
\begin{align}
    \boxed{{}^{x_0}R_{2a/0} = -\frac{I_{2a}}{r_{2a}^2}\,\ddot{x}^{(1)}
        + c_{2a,rr}\sigma\,{}^{z_0}R_{2a/0}
        + \frac{r_{2a,cg}}{r_{2a}}\,T_{2a/3b}
        - \frac{\sigma\,r_{2a,bk}}{r_{2a}}\,B_{2a/1}} \\
    \boxed{{}^{x_0}R_{2b/0} = -\frac{I_{2b}}{r_{2b}^2}\,\ddot{x}^{(1)}
        + c_{2b,rr}\sigma\,{}^{z_0}R_{2b/0}
        - \frac{\sigma\,r_{2b,bk}}{r_{2b}}\,B_{2b/1}}
\end{align}

\subsubsubsection{Normal reactions}
The rotation equation, after substituting kinematics and traction, yields the auxiliary quantity $K$ which collects all contributions to the pitching moment. The normal reactions follow from $K$ and the vertical equilibrium. Note that the EOM needs only the sum of the normal loads, so this step is post-processing:
\begin{align}
    \boxed{{}^{z_0}R_{2a/0} = \frac{K - l_D'\,Mg\cos\theta}{L}} \qquad
    \boxed{{}^{z_0}R_{2b/0} = -\frac{K + l_C'\,Mg\cos\theta}{L}}
\end{align}
\begin{equation}
    K = -\frac{J}{r_{2a}}\,\ddot{x}^{(1)}
        - hM\!\left(\ddot{x}^{(1)} + g\sin\theta\right)
        - r_{2a}\,A_{1,x} + A_{1,m},
    \qquad
    J = I_{2a} + I_{2b} + \frac{r_{2a,cg}}{r_{3b,cr}}\,I_{3b}
\end{equation}

\subsubsubsection{Equation of motion}
Substituting traction into the bike translation equation and using ${}^{z_0}R_{2a/0} + {}^{z_0}R_{2b/0} = -Mg\cos\theta$ collapses the rolling resistance terms. The effective mass $M_{eq}$ reflects all rotational inertias to the frame linear acceleration:
\begin{equation}
    \resizebox{\linewidth}{!}{$\displaystyle
    \boxed{
    M_{eq}\,\ddot{x}^{(1)} =
        \frac{r_{1,pedal}\,r_{2a,cg}}{r_{3b,cr}\,r_{2a}}\left(P_{3b/1,L} + P_{3b/1,R}\right)
        - \frac{\sigma}{r_{2a}}\left(r_{2a,bk}\,B_{2a/1} + r_{2b,bk}\,B_{2b/1}\right)
        - Mg\sin\theta
        - \sigma\,c_{2a,rr}\,Mg\cos\theta
        - A_{1,x}
    }$}
\end{equation}
\begin{equation}
    M_{eq} = M + \frac{I_{2a} + I_{2b}}{r_{2a}^2} + \frac{I_{3b}\,r_{2a,cg}^2}{r_{3b,cr}^2\,r_{2a}^2}
\end{equation}

\subsubsubsection{Limits}
The following conditions must hold for the model assumptions to remain valid. Violation of any one signals a regime change: wheel lift, chain slack, or sliding contact:
\begin{align}
    {}^{z_0}R_{2a/0} &\le 0, \qquad {}^{z_0}R_{2b/0} \le 0 \qquad \text{(no wheel lift)} \\
    \left|{}^{x_0}R_{2a/0}\right| &\le -\mu_s\,{}^{z_0}R_{2a/0}, \qquad
    \left|{}^{x_0}R_{2b/0}\right| \le -\mu_s\,{}^{z_0}R_{2b/0} \qquad \text{(no slip)} \\
    T_{2a/3b} &\ge 0 \qquad \text{(no freewheel)}
\end{align}

\subsection{Analysis}

\subsubsubsection{Assumptions}
\begin{itemize}
    \item Steady state --- $\ddot{x}^{(1)} = 0$, constant speed on constant gradient
    \item No braking --- $B_{2a/1} = B_{2b/1} = 0$
    \item Forward motion --- $\sigma = 1$
    \item Still air or head/tailwind only --- drag enters as $\tfrac{1}{2}\rho\,C_dA\,v_{rel}^2$ ($\rho$ here is air density)
    \item Drivetrain efficiency not modelled --- power input equals power at the chain
    \item All limits satisfied --- no slip, no wheel lift, no freewheel
\end{itemize}

\subsubsubsection{Equilibrium}
Setting $\ddot{x}^{(1)} = 0$, $B_{2a/1} = B_{2b/1} = 0$, $\sigma = 1$ gives the steady-state force balance. Three resistance terms appear: gradient, rolling, and aerodynamic drag, scaling respectively as $v^0$, $v^0$, and $v^2$:
\begin{equation}
    \frac{r_{1,pedal}\,r_{2a,cg}}{r_{3b,cr}\,r_{2a}}\left(P_{3b/1,L} + P_{3b/1,R}\right) = Mg\left(c_{2a,rr}\cos\theta + \sin\theta\right) + \tfrac{1}{2}\rho\,C_dA\,v^2
\end{equation}
Expressing the left side as prescribed power $\mathcal{P}$ using $|\dot{\phi}_{3b}| = r_{2a,cg}v/(r_{3b,cr}r_{2a})$ eliminates the gear geometry:
\begin{equation}
    \frac{\mathcal{P}}{v} = Mg\left(c_{2a,rr}\cos\theta + \sin\theta\right) + \tfrac{1}{2}\rho\,C_dA\,v^2
\end{equation}

\subsubsubsection{Sensitivity}
Write the aerodynamic drag power as $k \equiv \tfrac{1}{2}\rho\,C_dA$ and drop rolling resistance
(an order of magnitude below gravity and drag at any sustained speed), leaving the simplest model
that still contains both effects each question is about:
\begin{equation}
    \mathcal{P} = Mg\,v\sin\theta + k\,v^3
\end{equation}

\subparagraph{Elasticities.}
The elasticity $E_p \equiv (p/v)(\partial v/\partial p)$ gives the fractional change in speed per
fractional change in parameter $p$, at constant $\mathcal{P}$. Write the balance as
$G(v;\,M,k,\theta,\mathcal{P}) = Mg\,v\sin\theta + k\,v^3 - \mathcal{P} = 0$ and
differentiate implicitly:
\begin{equation}
    \frac{\partial v}{\partial p} = -\frac{\partial G/\partial p}{\partial G/\partial v},
    \qquad
    \frac{\partial G}{\partial v} = Mg\sin\theta + 3kv^2 \equiv D
\end{equation}
\begin{equation}
    \boxed{E_M = -\frac{Mg\sin\theta}{D}, \qquad
           E_k = -\frac{kv^2}{D}, \qquad
           E_{\mathcal{P}} = \frac{Mg\sin\theta + kv^2}{D}}
\end{equation}
Since $G$ is linear in $(M, k, \mathcal{P})$ at fixed $v$, Euler's theorem gives the exact identity:
\begin{equation}
    E_M + E_k + E_{\mathcal{P}} = 0
\end{equation}

\subparagraph{Two gradient limits.}
\textbf{Flat, aero-dominated} ($Mg\sin\theta \ll kv^2$, $D \approx 3kv^2$):
\begin{equation}
    E_M \approx 0, \qquad E_k \approx -\tfrac{1}{3}, \qquad E_{\mathcal{P}} \approx \tfrac{1}{3}
    \quad\Longrightarrow\quad v \propto \mathcal{P}^{1/3},\;\; v \propto k^{-1/3}
\end{equation}
Weight is irrelevant on the flat; speed depends only on power and drag, and weakly (cube root).

\textbf{Steep, gravity-dominated} ($kv^2 \ll Mg\sin\theta$, $D \approx Mg\sin\theta$):
\begin{equation}
    E_M \approx -1, \qquad E_k \approx 0, \qquad E_{\mathcal{P}} \approx 1
    \quad\Longrightarrow\quad v \approx \frac{\mathcal{P}}{Mg\sin\theta}
\end{equation}
Aerodynamics is irrelevant while climbing; speed is set entirely by power-to-weight.

\subparagraph{Does aerodynamics matter more for a lighter rider?}
Two riders with the same $\mathcal{P}$ and $k$, differing only in mass. Since $E_M < 0$ the
lighter rider is faster. Differentiate $|E_k| = kv^2/D$ along the equilibrium curve:
\begin{equation}
    \frac{\mathrm{d}|E_k|}{\mathrm{d}v} = \frac{2kv\,Mg\sin\theta}{D^2} > 0
\end{equation}
$|E_k|$ is strictly increasing in $v$, so aerodynamics matters strictly more to the lighter
(faster) rider. The gap is smallest while climbing ($|E_k|\approx 0$ for both) and largest on
the flat ($|E_k|\to\tfrac{1}{3}$ for both).

\subparagraph{Is there a best rider weight?}
Comparing riders of different size: power and frontal area both scale with mass. Model with
exponents against a reference rider $(M_0, \mathcal{P}_0, k_0)$:
\begin{equation}
    \mathcal{P}(M) = \mathcal{P}_0\!\left(\frac{M}{M_0}\right)^{a}, \qquad
    k(M) = k_0\!\left(\frac{M}{M_0}\right)^{b}
\end{equation}
($b \approx \tfrac{2}{3}$ from geometric similarity; $a$ is physiological, typically $a < 1$.)
The size elasticity follows by chain rule:
\begin{equation}
    \boxed{\left.\frac{\mathrm{d}\ln v}{\mathrm{d}\ln M}\right|_{\text{across riders}}
    = E_M + a\,E_{\mathcal{P}} + b\,E_k}
\end{equation}
Plugging in the two limits: climbing gives $\approx a - 1 < 0$ (smaller always wins for $a < 1$);
flat gives $\approx (a-b)/3$ (bigger wins if power scales faster than frontal area, $a > b$).
Best weight flips sign with gradient; the crossover depends on $a - b$, not on mass directly.

\subparagraph{Is a steeper climb more energy-efficient for the same height gain?}
Fix height $\Delta h$; route length is $d(\theta) = \Delta h / \sin\theta$. Energy spent at
constant power $\mathcal{P}$:
\begin{equation}
    E(\theta) = \mathcal{P}\,\frac{d(\theta)}{v(\theta)}
    = \frac{\mathcal{P}\,\Delta h}{v(\theta)\sin\theta}
\end{equation}
From the power balance, $v\sin\theta = (\mathcal{P} - kv^3)/(Mg)$, which is maximised by
minimising $v$. Since $\partial v/\partial\theta < 0$ (steeper is slower),
$v\sin\theta$ rises monotonically with gradient and $E(\theta)$ falls:
\begin{equation}
    \frac{\mathrm{d}E}{\mathrm{d}\theta} < 0
    \;\Longrightarrow\;
    \text{steeper is always more energy-efficient for a fixed height gain}
\end{equation}
Mechanism: a steeper route is shorter and slower; aerodynamic loss scales as $v^2 \times
\text{distance}$, so cutting speed helps cubically while the shorter distance helps linearly.

\subsubsubsection{Stoppie limit}
Front braking transfers load forward. The rear wheel lifts when ${}^{z_0}R_{2a/0} = 0$. Setting this in the normal reaction expression fixes $K$, which is then substituted into its definition to give the critical deceleration:
\begin{equation}
    {}^{z_0}R_{2a/0} = 0 \;\Rightarrow\; K = l_D'\,Mg\cos\theta
\end{equation}
\begin{equation}
    \ddot{x}^{(1)}_{\text{stoppie}} = \frac{-\left(l_D'\,Mg\cos\theta + hM\,g\sin\theta + r_{2a}\,A_{1,x} - A_{1,m}\right)}{J/r_{2a} + hM}
\end{equation}

\subsubsubsection{Wheelie limit}
Hard acceleration transfers load rearward. The front wheel lifts when ${}^{z_0}R_{2b/0} = 0$:
\begin{equation}
    {}^{z_0}R_{2b/0} = 0 \;\Rightarrow\; K = -l_C'\,Mg\cos\theta
\end{equation}
\begin{equation}
    \ddot{x}^{(1)}_{\text{wheelie}} = \frac{l_C'\,Mg\cos\theta - hM\,g\sin\theta - r_{2a}\,A_{1,x} + A_{1,m}}{J/r_{2a} + hM}
\end{equation}

\subsubsubsection{Traction limit}
Slip is checked after the normal reactions are known. It does not feed back into the EOM --- the no-slip assumption is either valid or it is not, and violation signals that a different model (e.g.\ with combined slip) is required. The available friction at each contact is $\mu_s\left|{}^{z_0}R_{k/0}\right|$; the longitudinal force must stay within it:
\begin{align}
    \left|{}^{x_0}R_{2a/0}\right| &\le -\mu_s\,{}^{z_0}R_{2a/0} \;\Rightarrow\; \text{rear contact valid} \\
    \left|{}^{x_0}R_{2b/0}\right| &\le -\mu_s\,{}^{z_0}R_{2b/0} \;\Rightarrow\; \text{front contact valid}
\end{align}

\subsubsubsection{Todo}
\begin{itemize}
    \item FBD / dimension diagrams
\end{itemize}
\end{aiblock}
